% =============================================================================
%  INTRODUCTION GÉNÉRALE
% =============================================================================
\cleardoublepage
\entete{Introduction générale}
\chapter*{INTRODUCTION GÉNÉRALE}
\addcontentsline{toc}{chapter}{INTRODUCTION GÉNÉRALE}

% -----------------------------------------------------------------------------
\section*{1. Contexte général de l'étude}
\addcontentsline{toc}{section}{1. Contexte général de l'étude}

À l'ère du numérique, le secteur culturel et patrimonial traverse une
transformation profonde à l'échelle mondiale. La valorisation numérique du
patrimoine — qu'il s'agisse de collections d'art, d'architectures
traditionnelles ou de savoir-faire artisanaux — est devenue un axe stratégique
majeur pour les institutions désireuses de préserver la mémoire collective tout
en élargissant leur audience. Longtemps réservée à des présentations physiques
et statiques, la médiation culturelle s'oriente désormais vers des expériences
interactives, portées par les avancées de l'ingénierie logicielle, de l'imagerie
en trois dimensions et de l'intelligence artificielle.

L'actualité récente a accéléré cette prise de conscience. Les crises sanitaires
mondiales et la montée en puissance du tourisme virtuel ont mis en évidence la
vulnérabilité des modèles reposant uniquement sur la présence physique des
visiteurs. Selon l'UNESCO (2021), près de 90~\% des musées du monde ont dû
fermer temporairement leurs portes durant la pandémie de COVID-19, et un musée
sur dix a craint de ne jamais rouvrir ; cette épreuve a poussé les institutions
à accélérer leur virage numérique pour maintenir le lien avec le public.
Aujourd'hui, les attentes des usagers ont évolué : le public recherche des accès
à distance personnalisés, immersifs et disponibles à tout moment. Dans le
contexte africain en particulier, où la richesse culturelle se heurte parfois à
des contraintes d'infrastructure, de mobilité ou d'éloignement géographique, le
numérique offre une occasion sans précédent de faire connaître le patrimoine
local à l'échelle internationale.

Un second phénomène pèse sur la manière dont ce patrimoine peut être présenté :
la dispersion. Une part considérable des objets africains conservés aujourd'hui
se trouve hors du continent. Le rapport remis à la présidence française par
Sarr et Savoy (2018) estime que 85 à 90~\% du patrimoine matériel africain se
trouverait hors d'Afrique. Concrètement, une chefferie ou une fondation
camerounaise peut posséder un objet dont les pièces apparentées — issues du même
atelier, de la même communauté ou de la même époque — sont conservées à des
milliers de kilomètres, dans des musées qui ne savent pas qu'elles ont une
famille. Ce constat déplace la question : il ne s'agit plus seulement de montrer
ce que l'on possède, mais aussi de rendre visible ce dont on a été séparé.

Malgré cet engouement mondial pour la numérisation culturelle, le constat sur le
terrain révèle d'importantes disparités. De nombreuses structures
socio-culturelles d'Afrique subsaharienne restent confrontées à des difficultés
d'accessibilité et de visibilité. La découverte de leurs espaces et de leurs
activités demeure strictement dépendante d'une visite sur place. Cette
contrainte engendre trois limites majeures :

\begin{itemize}
  \item un frein géographique et économique pour les publics éloignés,
        régionaux, nationaux ou internationaux, ainsi que pour la diaspora ;
  \item une barrière physique pour les personnes à mobilité réduite ;
  \item l'absence d'outils capables de renseigner le visiteur à distance de
        manière vivante et fiable.
\end{itemize}

C'est précisément dans ce cadre que s'inscrit la Fondation Jean Félicien Gacha,
implantée à Bangoulap, dans le département du Ndé, région de l'Ouest du
Cameroun. Véritable carrefour d'art, de culture, d'artisanat et de développement
communautaire, la Fondation abrite un espace d'une grande richesse
architecturale et patrimoniale : ateliers de création, espaces d'exposition,
résidences d'artistes, espaces verts. Elle héberge également le \emph{Think
Tank} Foudjem, cellule de réflexion et d'innovation au sein de laquelle notre
stage s'est déroulé du 20~juin au 20~août 2026. Consciente des enjeux de son
époque, la Fondation souhaite moderniser son offre et permettre une découverte
de ses lieux sans contrainte de déplacement. Actuellement, la valorisation de
ses richesses reste confinée à son site physique, ce qui constitue une limite
pour son rayonnement.

Pour répondre à ce besoin de modernisation et d'accessibilité à distance, trois
grands domaines technologiques convergent aujourd'hui :

\begin{itemize}
  \item \textbf{les systèmes d'information de gestion} : ils constituent la
        colonne vertébrale administrative de l'institution et le canal public
        par lequel ses services sont rendus accessibles ;
  \item \textbf{la modélisation en trois dimensions et les environnements
        virtuels} : ils apportent la dimension spatiale et visuelle, et
        permettent d'offrir une véritable expérience de déplacement à distance ;
  \item \textbf{l'intelligence artificielle conversationnelle} : elle dote
        l'environnement virtuel d'une capacité d'assistance, capable de répondre
        aux questions des visiteurs en s'appuyant sur les connaissances réelles
        de la Fondation.
\end{itemize}

Cette situation met en évidence un enjeu à la fois scientifique et pratique : la
simple mise en ligne de contenus ne suffit plus. Il s'agit de concevoir un
ensemble numérique cohérent, capable de combiner la gestion de l'institution,
l'immersion dans ses espaces et l'assistance au visiteur, pour offrir une
expérience de découverte fluide et instructive. C'est pour répondre à ce besoin
exprimé par la Fondation Jean Félicien Gacha que la présente étude a été menée.

% -----------------------------------------------------------------------------
\section*{2. Problématique de l'étude}
\addcontentsline{toc}{section}{2. Problématique de l'étude}

\subsection*{2.1. Présentation du problème}
\addcontentsline{toc}{subsection}{2.1. Présentation du problème}

À l'ère de la médiation culturelle numérique, la valorisation du patrimoine au
sein de la Fondation Jean Félicien Gacha — comme dans la majorité des
institutions culturelles de la sous-région — souffre d'un modèle dépassé : la
vitrine virtuelle statique. L'expérience proposée au public se réduit encore
trop souvent à une juxtaposition de photographies plates et de courtes fiches
textuelles figées.

Sur le terrain, cette façon de faire montre trois limites majeures.

\begin{enumerate}
  \item \textbf{Une visite passive, sans véritable sensation d'espace.} Une
        simple photographie ne permet ni de percevoir les volumes ni d'explorer
        les lieux. Le visiteur à distance ne peut pas se déplacer dans les
        espaces d'exposition de la Fondation, ni observer une œuvre sous tous
        ses angles, ni en apprécier la taille réelle. Il regarde une image ; il
        ne visite pas.

  \item \textbf{Des objets isolés et une histoire incomplète.} Les fiches
        descriptives traitent chaque pièce séparément. Les dispositifs actuels
        ne croisent pas leurs informations avec celles d'autres institutions
        pour indiquer si des œuvres apparentées — fabriquées par la même
        communauté, dans le même atelier ou à la même époque — existent
        ailleurs. Cet isolement empêche de mesurer la véritable portée
        historique des objets et prive le visiteur du récit qui les relie.

  \item \textbf{L'absence d'aide pour répondre aux questions précises.} Un texte
        descriptif court ne donne que des informations de base. Si un visiteur
        souhaite en savoir davantage sur un détail, une technique ou une
        histoire, il n'a aucun moyen d'interroger un guide capable de lui
        répondre immédiatement et, surtout, de lui répondre juste.
\end{enumerate}

À ces trois limites s'en ajoute une quatrième, d'ordre organisationnel. Les
solutions de visite virtuelle existantes sont généralement conçues pour une
institution unique. Or la Fondation n'est pas isolée : elle s'inscrit dans un
tissu de chefferies, de musées et d'associations patrimoniales qui rencontrent
le même besoin. Un dispositif taillé pour une seule structure devrait être
reconstruit intégralement pour la suivante, ce qui en interdit de fait la
diffusion.

Face à ces limites, l'enjeu ne consiste pas seulement à mettre des photographies
en ligne, mais à créer une plateforme numérique unifiée réunissant l'exploration
en trois dimensions, la mise en relation entre institutions et un assistant
capable de s'appuyer directement sur les documents de la Fondation.

\subsection*{2.2. Formulation du problème}
\addcontentsline{toc}{subsection}{2.2. Formulation du problème}

Le problème central de ce travail se résume à la question suivante.

\encadre{Question générale de l'étude}{%
Comment concevoir et réaliser une plateforme numérique unifiée capable d'offrir
une visite virtuelle en trois dimensions, de mettre en relation un réseau
d'institutions pour retrouver les œuvres apparentées, et d'assurer une
assistance conversationnelle fiable fondée sur les données réelles de la
Fondation Jean Félicien Gacha ?}

Pour répondre de manière claire à cette question, notre étude s'appuie sur
quatre questions spécifiques.

\begin{enumerate}
  \item[\textbf{QS1.}] Comment représenter les espaces et les œuvres de la
        Fondation pour remplacer les simples photographies par une exploration
        en trois dimensions fluide et interactive ?

  \item[\textbf{QS2.}] Comment relier les inventaires de la Fondation à ceux
        d'autres institutions afin d'identifier les pièces issues d'une même
        origine et d'enrichir ainsi leur valeur culturelle ?

  \item[\textbf{QS3.}] Comment concevoir un assistant conversationnel capable de
        comprendre les questions des visiteurs et d'y répondre exactement, en se
        fondant uniquement sur la documentation validée par la Fondation ?

  \item[\textbf{QS4.}] Comment construire l'architecture de l'application pour
        qu'il soit simple d'y accueillir de nouvelles institutions partenaires
        par la suite, sans tout reconstruire ?
\end{enumerate}

% -----------------------------------------------------------------------------
\section*{3. Hypothèses de l'étude}
\addcontentsline{toc}{section}{3. Hypothèses de l'étude}

\subsection*{3.1. Hypothèse générale}
\addcontentsline{toc}{subsection}{3.1. Hypothèse générale}

\encadre{Hypothèse générale}{%
La mise en place, au sein de la Fondation Jean Félicien Gacha, d'une plateforme
numérique unifiée réunissant une visite à 360 degrés, la modélisation en trois
dimensions des œuvres, un assistant conversationnel ancré sur la documentation
de l'institution et une architecture ouverte à plusieurs organisations, permet
de restituer à distance l'essentiel de l'expérience de visite et d'élargir
l'audience de la Fondation au-delà de ses murs.}

\subsection*{3.2. Hypothèses spécifiques}
\addcontentsline{toc}{subsection}{3.2. Hypothèses spécifiques}

\begin{enumerate}
  \item[\textbf{HS1.}] La restitution des espaces par des panoramas à 360 degrés
        reliés entre eux, complétée par la présentation des œuvres en trois
        dimensions et en réalité augmentée, procure au visiteur distant une
        perception des volumes et des dimensions que la photographie plate ne
        permet pas.

  \item[\textbf{HS2.}] Le rapprochement automatique des notices de la Fondation
        avec celles des collections ouvertes accessibles en ligne permet
        d'identifier des œuvres apparentées et d'enrichir la valeur documentaire
        des pièces conservées à Bangoulap.

  \item[\textbf{HS3.}] Un assistant conversationnel dont les réponses sont
        construites à partir des seuls documents validés par la Fondation
        fournit des réponses plus exactes et plus vérifiables qu'un assistant
        répondant de mémoire, et réduit le risque d'affirmations inventées.

  \item[\textbf{HS4.}] Une architecture qui isole les données de chaque
        institution dès la base de données permet d'accueillir une nouvelle
        organisation sans modifier le code de l'application ni compromettre la
        confidentialité des données des autres.
\end{enumerate}

% -----------------------------------------------------------------------------
\section*{4. Objectifs de l'étude}
\addcontentsline{toc}{section}{4. Objectifs de l'étude}

\subsection*{4.1. Objectif général}
\addcontentsline{toc}{subsection}{4.1. Objectif général}

L'objectif général de cette étude est le suivant :

\encadre{Objectif général}{%
Concevoir et réaliser, pour la Fondation Jean Félicien Gacha, une plateforme
numérique qui permette à toute personne, où qu'elle se trouve, de visiter la
Fondation, de découvrir ses œuvres et d'obtenir des réponses à ses questions,
sans avoir à se déplacer.}

Dit simplement : il s'agit de faire entrer la Fondation dans l'écran de ceux qui
ne peuvent pas franchir son portail. Un visiteur installé à Douala, à Paris ou à
Montréal doit pouvoir se promener dans les salles, s'approcher d'un objet, en
faire le tour, savoir à quoi il servait et qui l'a possédé, et poser une
question comme il le ferait à un guide présent à ses côtés.

\subsection*{4.2. Objectifs spécifiques}
\addcontentsline{toc}{subsection}{4.2. Objectifs spécifiques}

Quatre objectifs spécifiques découlent des quatre questions posées plus haut.

\begin{enumerate}
  \item[\textbf{OS1.}] \textbf{Rendre la visite possible à distance.} Reproduire
        les espaces de la Fondation sous forme de vues à 360 degrés reliées
        entre elles, dans lesquelles le visiteur se déplace de salle en salle,
        et présenter les œuvres en volume afin qu'il puisse les examiner sous
        tous les angles et se rendre compte de leur taille réelle.

  \item[\textbf{OS2.}] \textbf{Replacer chaque œuvre dans sa famille.} Mettre en
        place un dispositif qui recherche, dans les catalogues ouverts d'autres
        musées, les objets proches de ceux de la Fondation, présente les
        rapprochements trouvés avec leur justification, et laisse au
        conservateur le soin de les valider ou de les écarter.

  \item[\textbf{OS3.}] \textbf{Donner une voix à la Fondation.} Concevoir un
        assistant qui répond aux questions des visiteurs en s'appuyant
        uniquement sur les textes fournis par l'institution, qui cite ce sur
        quoi il se fonde, et qui reconnaît son ignorance plutôt que d'inventer
        une réponse.

  \item[\textbf{OS4.}] \textbf{Ouvrir la porte aux autres institutions.}
        Construire l'application de telle sorte qu'une nouvelle chefferie, un
        nouveau musée ou une nouvelle fondation puisse s'y inscrire, obtenir son
        propre espace et ses propres adresses, et gérer ses collections sans
        jamais voir celles des autres.
\end{enumerate}

% -----------------------------------------------------------------------------
\section*{5. Justification de l'étude}
\addcontentsline{toc}{section}{5. Justification de l'étude}

\subsection*{5.1. Au plan scientifique}
\addcontentsline{toc}{subsection}{5.1. Au plan scientifique}

Sur le plan scientifique, ce travail apporte une contribution à trois niveaux.

D'abord, il documente une question encore peu traitée dans la littérature
francophone : celle de la visite immersive appliquée aux institutions
patrimoniales d'Afrique subsaharienne, c'est-à-dire à des structures qui
disposent d'un patrimoine considérable mais de moyens techniques limités. Les
travaux existants portent majoritairement sur de grands musées occidentaux,
dotés de budgets et d'équipes sans commune mesure. Décrire ce qui fonctionne —
et surtout ce qui ne fonctionne pas — dans un contexte de ressources contraintes
constitue en soi un apport.

Ensuite, l'étude produit une mesure comparative rarement publiée : celle du gain
apporté par l'enrichissement automatique d'une requête de recherche
documentaire. Nous avons observé qu'une recherche formulée en français dans les
catalogues internationaux ramène beaucoup de résultats de faible pertinence,
tandis que la même recherche, reformulée dans le vocabulaire de catalogage
employé par ces institutions, en ramène moins mais de bien meilleure qualité.
Ce résultat, chiffré au chapitre~3, montre que l'échec initial d'une recherche
ne traduit pas nécessairement une absence de données, mais souvent un décalage
de vocabulaire.

Enfin, le travail met à l'épreuve, sur un cas réel, le principe de la génération
augmentée par la recherche appliquée à un domaine où l'exactitude prime sur la
fluidité. En médiation culturelle, une réponse inventée n'est pas une
maladresse : c'est une falsification. L'étude propose et évalue des garde-fous
concrets pour éviter cette dérive.

\subsection*{5.2. Au plan pratique}
\addcontentsline{toc}{subsection}{5.2. Au plan pratique}

Sur le plan pratique, l'utilité du travail est immédiate et se laisse résumer en
quelques phrases simples.

\begin{itemize}
  \item \textbf{Pour la Fondation}, la plateforme est un outil de travail. Le
        personnel dispose enfin d'un endroit unique où sont enregistrés les
        musées, les salles, les œuvres, les tarifs, les événements, les
        commandes et les visiteurs. Ce qui se faisait sur des cahiers et des
        fichiers dispersés se fait désormais dans un seul système, consultable
        et sauvegardé.

  \item \textbf{Pour le public éloigné}, elle supprime la contrainte du
        déplacement. Un membre de la diaspora, un chercheur étranger, un élève
        d'une école située à plusieurs centaines de kilomètres, une personne à
        mobilité réduite : tous peuvent visiter.

  \item \textbf{Pour le rayonnement de l'institution}, elle transforme une
        visite en relation durable. Le visiteur qui crée un compte devient un
        adhérent que la Fondation peut informer de ses expositions et de ses
        événements, ce qui n'était pas possible auparavant.

  \item \textbf{Pour les autres institutions de la région}, elle offre un modèle
        reproductible. L'architecture retenue permet à une chefferie voisine de
        s'inscrire et de disposer de son propre espace en quelques minutes,
        sans développement supplémentaire. Le coût de la mise en ligne du
        patrimoine s'en trouve considérablement réduit.

  \item \textbf{Pour l'économie locale}, enfin, la boutique et la billetterie
        intégrées ouvrent un canal de vente aux artisans et aux ateliers de la
        Fondation, dont la production restait jusqu'ici tributaire des visites
        physiques.
\end{itemize}

% -----------------------------------------------------------------------------
\section*{6. Délimitation de l'étude}
\addcontentsline{toc}{section}{6. Délimitation de l'étude}

Conformément aux recommandations méthodologiques, la présente étude est délimitée
sous cinq angles.

\begin{itemize}
  \item \textbf{Délimitation spatiale.} L'étude porte sur la Fondation Jean
        Félicien Gacha, sise à Bangoulap, département du Ndé, région de l'Ouest
        du Cameroun, et plus précisément sur son \emph{Think Tank} Foudjem.
        Les données ont été recueillies sur ce site.

  \item \textbf{Délimitation temporelle.} Le travail de terrain s'est déroulé du
        20~juin au 20~août 2026, soit une période de deux mois correspondant à
        la durée du stage académique. Les mesures de performance et de coût
        rapportées portent sur cette même période.

  \item \textbf{Délimitation thématique.} L'étude traite de la restitution
        numérique de la visite : représentation des espaces, présentation des
        œuvres, assistance au visiteur et mise en relation documentaire. Elle
        n'aborde ni la conservation matérielle des objets, ni la muséographie
        physique, ni les questions juridiques de restitution des biens
        culturels, qui relèvent d'autres disciplines.

  \item \textbf{Délimitation conceptuelle.} Par \emph{visite immersive}, nous
        entendons ici un parcours numérique dans lequel l'utilisateur change de
        point de vue et de lieu par ses propres actions ; nous excluons les
        dispositifs à casque de réalité virtuelle, dont le coût et l'équipement
        les rendent inaccessibles au public visé. Par \emph{réalité augmentée},
        nous entendons la projection d'un objet à sa taille réelle dans
        l'environnement de l'utilisateur au moyen de la caméra de son téléphone.

  \item \textbf{Délimitation de la population.} L'étude s'adresse à trois
        catégories d'acteurs : le personnel de la Fondation, les visiteurs
        présents sur le site et les visiteurs à distance. Les mineurs non
        accompagnés ont été exclus de l'enquête.

  \item \textbf{Portée et limites de validité.} S'agissant d'une étude de cas
        conduite auprès d'une institution unique et sur un échantillon
        volontairement réduit, les résultats ne prétendent pas à une
        généralisation statistique. Ils décrivent un dispositif, en mesurent le
        comportement et en tirent des enseignements transférables, sous réserve
        de vérification, aux institutions présentant des caractéristiques
        comparables.
\end{itemize}

% -----------------------------------------------------------------------------
\section*{7. Plan du rapport}
\addcontentsline{toc}{section}{7. Plan du rapport}

Hormis l'introduction générale et la conclusion générale, le présent rapport
s'organise en quatre chapitres.

Le \textbf{chapitre~1} établit le cadre conceptuel et l'état de l'art : il
définit les notions mobilisées et examine les solutions existantes ainsi que
leurs limites. Le \textbf{chapitre~2} expose la méthodologie : nature de
l'étude, variables et indicateurs, population et échantillon, outils de
collecte, techniques d'analyse et méthode d'ingénierie logicielle retenue.

Le \textbf{chapitre~3} présente le site de l'étude, le déroulement du stage, les
données recueillies et les résultats obtenus, fonctionnalité par fonctionnalité.
Le \textbf{chapitre~4} procède à l'analyse et au diagnostic de la situation,
vérifie les hypothèses de départ et présente l'intervention proposée, son coût,
sa faisabilité et ses conditions de pérennité.

Le rapport s'achève sur une conclusion générale, un ensemble de perspectives,
les références bibliographiques et les annexes.
